%% ============================================================
%%  SUPPLEMENTARY MATERIAL A
%%  Routing Improvements for Hybrid LLM–Neural-Network Systems
%%  Impact on Extrapolation Performance in the DeFi Benchmark
%%
%%  Companion to: jmlr-hypatiax-paper-final.tex
%%  Rewritten & restructured: March 2026
%% ============================================================

\documentclass[11pt,a4paper]{article}

\usepackage[T1]{fontenc}
\usepackage[utf8]{inputenc}
\usepackage{geometry}
\geometry{margin=2.5cm}

\usepackage{amsmath,amssymb,amsthm}
\usepackage{booktabs}
\usepackage{array}
\usepackage{xcolor}
\usepackage{listings}
\usepackage{caption}
\usepackage{subcaption}
\usepackage{hyperref}
\usepackage{enumitem}
\usepackage{titlesec}
\usepackage{parskip}

%% ── Code listings ────────────────────────────────────────────────────────────
\definecolor{codeblue}{RGB}{0,90,170}
\definecolor{codegray}{RGB}{100,100,100}
\definecolor{codegreen}{RGB}{0,128,64}
\definecolor{codered}{RGB}{180,30,30}
\definecolor{backgray}{RGB}{248,248,248}

\lstset{
  language=Python,
  basicstyle=\ttfamily\small,
  keywordstyle=\color{codeblue}\bfseries,
  commentstyle=\color{codegray}\itshape,
  stringstyle=\color{codegreen},
  backgroundcolor=\color{backgray},
  frame=single,
  framerule=0.4pt,
  rulecolor=\color{black!20},
  breaklines=true,
  breakatwhitespace=true,
  numbers=left,
  numberstyle=\tiny\color{codegray},
  numbersep=8pt,
  xleftmargin=12pt,
  tabsize=4,
  showstringspaces=false,
  captionpos=b
}

\hypersetup{colorlinks=true, linkcolor=codeblue, citecolor=codeblue, urlcolor=codeblue}

\newtheorem{observation}{Observation}
\newtheorem{remark}{Remark}

%% ============================================================
\begin{document}

\title{%
  \textbf{Supplementary Material A}\\[6pt]
  \textbf{Routing Improvements for the Hybrid LLM--Neural-Network System}\\[4pt]
  \large Impact on DeFi Extrapolation Benchmark Performance
}
\author{Technical Note --- Hybrid System Development Log}
\date{February--March 2026}
\maketitle
\tableofcontents
\bigskip

\noindent\textbf{Note to reader.}  This document is Supplementary Material~A to
the main JMLR submission \emph{HypatiaX: A Hybrid Symbolic-Neural Framework for
Extrapolation-Reliable Analytical Discovery} (\texttt{jmlr-hypatiax-paper-final.tex}).
It provides full implementation details for the six routing fixes summarised in Section~7.4 (\emph{Component 3: Five-Stage Routing and Ensembling}) of the main paper---which now includes Proposition~1 (Routing Cascade Convergence), a formal guarantee that Phase~2 warm-starting is monotonically non-degrading over the Pareto front---and reports cumulative impact on the DeFi
extrapolation benchmark.

%% ============================================================
\section{Background and Motivation}
\label{sec:background}

\begin{remark}[Development-set status]
The DeFi equations used to motivate and validate the routing fixes in this
document (Section~\ref{sec:fixes}) overlap with equations in the Core~15
benchmark and the DeFi~74-case extrapolation suite (v3.0).  Post-fix performance
figures on the DeFi suite should therefore be interpreted as
\emph{in-sample routing improvement}, not held-out generalisation.
This is noted explicitly in Remark~3.2 of the main paper.
\end{remark}

The DeFi extrapolation benchmark evaluates three methods --- \emph{Pure LLM},
\emph{Neural Network (NN)}, and \emph{Hybrid} --- across 74 test cases ranging
in difficulty from easy linear formulas to hard transcendental functions.
Each case splits data so the test set lies \emph{outside} the training feature
range, directly probing extrapolation ability rather than interpolation.

Prior to the improvements described here, the aggregate picture was:

\begin{table}[h]
\centering
\caption{Baseline performance before routing improvements (standard cases,
         intractable cases excluded).}
\label{tab:baseline}
\begin{tabular}{lcccc}
\toprule
Method & Median $R^2$ & $R^2 > 0.9$ (\%) & $R^2 > 0.99$ (\%) & Catastrophic \\
\midrule
Pure LLM    & 1.000 & 85.7 & 83.9 & 4 \\
Neural Net  & 0.752 & 27.9 &  1.5 & 6 \\
Hybrid      & 1.000 & 73.5 & 66.2 & 1 \\
\bottomrule
\end{tabular}
\end{table}

Despite matching the Pure LLM's median $R^2 = 1.000$, the Hybrid's
\emph{tail performance} was substantially worse: 66.2\% vs.\ 83.9\% above
$R^2 > 0.99$, an 18\,pp gap.  Post-mortem analysis identified the root cause
as \textbf{routing confidence}: the hybrid's decision to assign a case to the
NN or blended ensemble relied on in-distribution training $R^2$ as the sole
criterion, which is a poor proxy for extrapolation quality.

Two additional data-generation bugs contributed:

\begin{itemize}[leftmargin=*]
  \item \textbf{Reserve Ratio / Spot Price (zero-variance targets).}
    Both cases generated \texttt{reserve\_x} and \texttt{reserve\_y} from
    identical \texttt{linspace} ranges, producing a ratio $\approx 1.0$ with
    negligible variance.
  \item \textbf{Impermanent Loss Breakeven (misclassified as tractable).}
    This case produced catastrophic NN scores ($R^2 \approx -41\,000$) and a
    degenerate LLM artefact ($R^2 = 1.000$, a train/test split artefact).
\end{itemize}

%% ============================================================
\section{Data-Generation Fixes}
\label{sec:data_fixes}

\subsection{Reserve Ratio and Spot Price (Fix~0)}

\begin{lstlisting}[caption={Buggy generator — identical linspace ranges.}]
reserve_x = np.linspace(100, 100000, n)
reserve_y = np.linspace(100, 100000, n)   # same range!
reserve_ratio = reserve_x / (reserve_y + 1e-10)  # ~1.0 everywhere
\end{lstlisting}

\begin{lstlisting}[caption={Fixed generator — independent log-uniform sampling (seed=7).}]
rng = np.random.default_rng(7)
reserve_x = np.exp(rng.uniform(np.log(100), np.log(100000), n))
reserve_y = np.exp(rng.uniform(np.log(100), np.log(100000), n))
reserve_ratio = reserve_x / reserve_y
# std(reserve_ratio) ~ 26.4,  range: [0.002, 212.6]
\end{lstlisting}

The same fix applies to the Spot Price generator
(\texttt{reserve\_a}, \texttt{reserve\_b}).  After the fix, both cases present
a genuine regression target.

\subsection{Impermanent Loss Breakeven (Fix~0b)}

\begin{lstlisting}[caption={Intractable flag added to IL Breakeven catalogue entry.}]
{
  "name": "Impermanent loss breakeven",
  "domain": "liquidity",
  "extrapolation_intractable": True,
  "intractable_reason":
    "Breakeven condition involves implicit solving; aggressive high-split "
    "puts test set in a degenerate regime. NN R2 collapses catastrophically; "
    "LLM 1.000 is a split artefact."
}
\end{lstlisting}

Excluding this case from the standard aggregate removes the last remaining
catastrophic case from the Hybrid column and avoids inflating the LLM
advantage.

%% ============================================================
\section{Routing Improvements (Fixes~1--4)}
\label{sec:routing_fixes}

All four routing improvements are implemented as methods of the
\texttt{EnhancedExtrapolationTest} class and fire \emph{after} the hybrid's
own training-time decision, potentially overriding it before test-set
evaluation.  The override priority order is:

\begin{enumerate}
  \item Fix~2 (formula complexity) — cheapest, no extra training.
  \item Fix~1 (extrapolation probe) — if Fix~2 did not already override.
  \item Fix~3 (LLM feature augmentation) — when decision remains \texttt{"nn"}.
  \item Fix~4 (distance-gated blend weight) — in \texttt{"ensemble"} paths.
\end{enumerate}

\subsection{Fix~1: Extrapolation Probe Degradation}

\paragraph{Motivation.}
A neural network fitting training data well may still extrapolate poorly.
In-distribution $R^2$ cannot detect this.

\paragraph{Method.}

\begin{enumerate}
  \item Sort training samples by the primary feature.
  \item Reserve the top $p = 15\%$ as a \emph{probe set}.
  \item Train a small MLP (layers: $[64, 32]$, 200 epochs, Adam) on the
        remaining $1-p$ fraction.
  \item Compute $\Delta R^2 = R^2_{\text{in}} - R^2_{\text{probe}}$.
        If $\Delta R^2 \geq 0.15$, override to \texttt{"llm"}.
\end{enumerate}

\begin{lstlisting}[caption={Extrapolation probe core logic.}]
def _extrapolation_probe_degradation(self, X_train, y_train, p=0.15):
    n   = len(X_train)
    split = int(n * (1 - p))
    order = np.argsort(X_train[:, 0])
    X_sorted, y_sorted = X_train[order], y_train[order]

    X_fit,   y_fit   = X_sorted[:split],  y_sorted[:split]
    X_probe, y_probe = X_sorted[split:],  y_sorted[split:]

    # ... train MLP on (X_fit, y_fit) ...

    degradation = max(0.0, r2_in - r2_probe)
    return degradation   # caller: if degradation >= 0.15 -> route to LLM
\end{lstlisting}

\paragraph{Measured gain.} $+6$ percentage points on $R^2 > 0.99$.

\subsection{Fix~2: Formula Complexity Routing}

\paragraph{Motivation.}
The LLM's extracted formula is free information about the functional class of
the target.  Transcendental or piecewise operations make NN extrapolation
structurally unreliable regardless of training $R^2$.

\paragraph{Method.}

\begin{lstlisting}[caption={Transcendental token list.}]
_TRANSCENDENTAL_TOKENS = [
    "math.exp", "np.exp", "exp(",
    "math.log", "np.log", "log(",
    "math.sqrt", "np.sqrt", "sqrt(",
    "norm.cdf", "norm.pdf", "scipy.stats",
    "np.maximum", "np.minimum", "max(", "min(",
    "math.sin", "np.sin", "math.cos", "np.cos",
    "**0.",   # fractional powers
]
\end{lstlisting}

Any match forces an immediate override to \texttt{"llm"} without running the
probe (cheaper, deterministic).  Fix~2 is checked before Fix~1.

\paragraph{Measured gain.} $+5$ percentage points.

\subsection{Fix~3: LLM Predictions as a Neural-Network Feature}

\paragraph{Motivation.}
When routing remains \texttt{"nn"}, the NN must infer the functional form from
scratch.  Concatenating LLM predictions as an extra input feature forces the
NN to learn residuals rather than the full function.

\paragraph{Method.}

\[
  \mathbf{X}_{\text{aug}}
  = \bigl[\mathbf{X} \;\big|\; \hat{y}_{\text{LLM}}\bigr]
  \in \mathbb{R}^{n \times (d+1)}.
\]

\begin{lstlisting}[caption={LLM-feature augmentation.}]
llm_pred_train = execute_python_code_get_predictions(llm_code, X_train)
llm_pred_test  = execute_python_code_get_predictions(llm_code, X_test)

X_train_aug = np.column_stack([X_train, llm_pred_train])
X_test_aug  = np.column_stack([X_test,  llm_pred_test])

nn_m = self._train_and_eval_nn(X_train_aug, y_train, X_test_aug, y_test)
\end{lstlisting}

Guard condition: requires positive LLM training $R^2$; otherwise plain NN.

\paragraph{Measured gain.} $+1$ percentage point.

\subsection{Fix~4: Distance-Gated Blend Weights}

\paragraph{Motivation.}
In the \texttt{"ensemble"} path, the LLM--NN blend weight was fixed at training
time.  In an extrapolation benchmark the test set is always outside the
training range by construction.

\paragraph{Method.}
Define the out-of-range fraction:
\[
  \rho = \frac{1}{m}\sum_{i=1}^{m}
         \mathbf{1}\!\left[
           \exists\,j : x_{ij} < \min_k x_{kj}^{\text{train}}
                       \;\text{or}\;
                       x_{ij} > \max_k x_{kj}^{\text{train}}
         \right].
\]
The LLM blend weight becomes:
\[
  w_{\text{LLM}} = w_0 + (1-w_0)\,\rho, \qquad w_0 = 0.3,
\]
so $w_{\text{LLM}} \in [0.3,\,1.0]$.  Since all extrapolation test sets satisfy
$\rho \approx 1$, the blend collapses toward pure LLM.

\begin{lstlisting}[caption={Distance-gated blend weight.}]
def _distance_llm_weight(X_test, X_train, base_weight=0.3):
    train_min = X_train.min(axis=0)
    train_max = X_train.max(axis=0)
    outside   = np.mean(
        np.any((X_test < train_min) | (X_test > train_max), axis=1)
    )
    return float(base_weight + (1.0 - base_weight) * outside)
\end{lstlisting}

\paragraph{Measured gain.} $+1$ percentage point.

%% ============================================================
\section{Fix~5: Unified Formula Evaluator and Routing Guard}
\label{sec:fix5}

\subsection{Root Cause: Divergent Evaluation Paths}

Two cases returned catastrophic Hybrid scores even though routing was correct
and the LLM formula was valid:

\begin{itemize}[leftmargin=*]
  \item \textbf{Correlated Portfolio VaR:} Pure LLM $= +1.000$,
    Hybrid $= -12.121$.
  \item \textbf{Impermanent loss in constant product:} Pure LLM $= +1.000$,
    Hybrid $= \text{NaN}$.
\end{itemize}

The Hybrid's test-set path called \texttt{hybrid.evaluate\_llm\_formula()},
while the Pure LLM path called
\texttt{PureLLMBaseline.test\_formula\_accuracy()}.  On formulas involving
matrix operations or multi-variable lookups, the two namespace implementations
diverged silently.

\subsection{Fix~5: Unified Evaluator}

\begin{lstlisting}[caption={Unified formula evaluator (Fix~5).}]
@staticmethod
def _eval_formula_r2(llm_code, X, y_true, var_names):
    try:
        from hypatiax.experiments.tests.hybrid_ensemble_system_defi_domain import (
            execute_python_code_get_predictions,
        )
        y_pred = execute_python_code_get_predictions(llm_code, X)
        if y_pred is None or np.any(np.isnan(y_pred)) or np.any(np.isinf(y_pred)):
            return float("nan"), False
        ss_res = np.sum((y_true - y_pred) ** 2)
        ss_tot = np.sum((y_true - np.mean(y_true)) ** 2)
        r2 = float(1 - ss_res / ss_tot) if ss_tot > 1e-10 else 0.0
        return r2, True
    except Exception:
        return float("nan"), False
\end{lstlisting}

\subsection{Routing Override Guard (Fix~5b)}

Two further cases exposed a related failure: Fixes~1 and~2 overrode to
\texttt{"llm"} even when the LLM formula had negative training $R^2$.

\begin{lstlisting}[caption={LLM-formula trustworthiness guard.}]
_llm_train_r2 = float(
    train_result.get("llm_result", {})
    .get("metrics", {}).get("r2", 0.0) or 0.0
)
_llm_formula_trustworthy = has_formula and (_llm_train_r2 > 0.0)

# Fixes 1 and 2 only fire when the formula actually fit training data
if _llm_formula_trustworthy and decision in ("nn", "ensemble"):
    if self._formula_has_transcendental(llm_code):   # Fix 2
        decision = "llm"

if _llm_formula_trustworthy and decision in ("nn", "ensemble"):
    if degradation >= 0.15:                           # Fix 1
        decision = "llm"
\end{lstlisting}

The same guard applies to Fix~3: if the LLM formula has negative training
$R^2$, augmentation is skipped.

\subsection{Case-Level Impact of Fix~5}

\begin{table}[h]
\centering
\caption{Case-level outcomes resolved by Fix~5 and routing guard.}
\label{tab:fix5_cases}
\begin{tabular}{llrrl}
\toprule
\textbf{Case} & \textbf{Type} & \textbf{LLM} & \textbf{Hybrid (before)} & \textbf{Hybrid (after)} \\
\midrule
Correlated Portfolio VaR & quadratic     & $+1.000$ & $-12.121$ & $\approx+1.000$ \\
IL in constant product   & alg.\ + sqrt  & $+1.000$ & NaN       & $\approx+1.000$ \\
Capital efficiency       & rational      & NaN      & $-3274$   & $\approx+0.827$ \\
Concentrated liquidity   & alg.\ + sqrt  & $-20.8$  & $-1606$   & $\approx+0.963$ \\
\bottomrule
\end{tabular}
\end{table}

Net effect: catastrophic count $3 \to 1$ (only \emph{Component ES} remains,
where all three methods fail equally).

%% ============================================================
\section{Projected Impact on the Leaderboard}
\label{sec:leaderboard}

\begin{table}[h]
\centering
\caption{Measured and projected $R^2 > 0.99$ (\%) by fix.
         Denominator $= 66$ standard cases (pre-v3.0 historical). The v3.0 benchmark uses $n=74$; confirmed final figure: \textbf{89.2\%} (see main paper Table~1).}
\label{tab:projected}
\begin{tabular}{llccc}
\toprule
\textbf{Fix} & \textbf{Description} & \textbf{Gain (pp)} & \textbf{Cumulative (\%)} & \textbf{Status} \\
\midrule
Baseline    & In-dist. $R^2$ routing (old)      & ---  & 66.2 & historic  \\
Fix~0       & Reserve Ratio / Spot Price        & $+2$ & 68.2 & measured  \\
Fix~0b      & IL Breakeven flagged intractable  & $+1$ & 69.2 & measured  \\
Fix~1       & Extrapolation probe routing       & $+6$ & 75.2 & measured  \\
Fix~2       & Formula complexity routing        & $+5$ & 80.2 & measured  \\
Fix~3       & LLM feature augmentation          & $+1$ & 81.2 & measured  \\
Fix~4       & Distance-gated blend weight       & $+1$ & 82.2 & measured  \\
\midrule
\multicolumn{2}{l}{\textbf{First complete run} (honest $n{=}66$)} & & \textbf{72.7} & measured \\
\multicolumn{2}{l}{Pure LLM baseline (honest $n{=}66$)}           & & 69.7          & measured \\
\midrule
Fix~5 (proj.) & Unified evaluator + routing guard & $+3$ & \textbf{75.8} & projected \\
\midrule
\textbf{v3.0 confirmed} & All fixes + benchmark corrections & --- & \textbf{89.2} & \textbf{measured} \\
\bottomrule
\end{tabular}
\end{table}

\begin{observation}
The first complete benchmark run confirms the Hybrid already
\emph{beats} the Pure LLM on the honest (NaN-penalty) metric:
$72.7\%$ vs $69.7\%$, a net advantage of $+3$ percentage points.
In the v3.0 benchmark ($n=74$), the confirmed final figure is
\textbf{89.2\%} near-perfect success rate ($R^2>0.99$)---within
the projected 88--90\% range---with zero catastrophic failures.
See main paper Table~1 for full verified figures.
\end{observation}

\begin{remark}
The raw benchmark output (83.6\% LLM, 77.4\% Hybrid) uses different
denominators per method (NaN excluded per method) and is not comparable
across methods.  All comparisons in this document use the fixed denominator
of 66.
\end{remark}

%% ============================================================
\section{Interaction Between Fixes}
\label{sec:interactions}

\begin{itemize}[leftmargin=*]
  \item \textbf{Fix~2 subsumes a subset of Fix~1.}  Cases with transcendental
    tokens are caught by Fix~2 before Fix~1 runs, avoiding the probe NN cost.
  \item \textbf{Fix~3 benefits from Fix~1's precision.}  Cases remaining
    \texttt{"nn"} after the probe are safer for LLM-feature augmentation.
  \item \textbf{Fix~4 becomes redundant for cases upgraded by Fixes~1/2.}
    Fix~4's value lies in residual ensemble cases that neither Fix~1 nor
    Fix~2 catches.
  \item \textbf{Data Fix~0 is orthogonal to routing.}  Its contribution
    is purely additive.
\end{itemize}

The gap between the per-fix sum (82.2\%) and the first-complete-run result
(72.7\%) reflects: (i) fix interactions where Fix~2 subsumes Fix~1 cases;
(ii) the shift from the NaN-excluded denominator used in per-fix measurement
to the honest $n{=}66$ fixed denominator.

%% ============================================================
\section{Remaining Open Issues}
\label{sec:open_issues}

\begin{enumerate}[leftmargin=*]
  \item \textbf{Formula evaluator divergence.}  Closed by Fix~5 — all formula
    evaluation now routes through a single code path.
  \item \textbf{Stochastic NN initialisation.}  Cases routed to \texttt{"nn"}
    use a freshly trained MLP, introducing variance.  A deterministic seed or
    model caching would eliminate this.
  \item \textbf{Probe threshold sensitivity.}  The $\Delta R^2 \geq 0.15$
    threshold was chosen heuristically.  A systematic sweep over
    $\{0.05, 0.10, 0.15, 0.20, 0.25\}$ on a held-out validation set would
    yield a more principled value.
  \item \textbf{Multi-feature probe.}  The probe currently sorts by the
    primary feature only.  For cases with multiple features of similar
    importance, a Mahalanobis-distance-based probe would be more robust.
  \item \textbf{LLM API timeouts.}  Two scores on cases~6--7 remain NaN
    due to connection errors.  These should be retried with exponential
    backoff ($\leq 3$ attempts, delays 5/15/45\,s).
\end{enumerate}

%% ============================================================
\section{Summary of Changes}
\label{sec:changes}

\begin{table}[h]
\centering
\caption{All changes applied across the two modified files.}
\label{tab:changes}
\begin{tabular}{llll}
\toprule
\textbf{File} & \textbf{Fix} & \textbf{Location} & \textbf{Description} \\
\midrule
\texttt{experiment\_protocol\_defi.py}
  & 0  & Test~5  & Reserve Ratio: independent log-uniform sampling \\
  & 0  & Test~7  & Spot Price: independent log-uniform sampling \\
\midrule
\texttt{test\_enhanced\_defi\_extrapolation.py}
  & 0b & Catalogue     & IL Breakeven flagged intractable \\
  & -- & \texttt{\_load\_checkpoint} & Deduplication by test-case name \\
  & 1  & \texttt{\_extrapolation\_probe\_degradation} & Probe method added \\
  & 2  & \texttt{\_formula\_has\_transcendental} & Token list + check added \\
  & 1,2& \texttt{test\_method} (hybrid block) & Override logic injected \\
  & 3  & \texttt{test\_method} (nn branch)     & LLM-feature augmentation \\
  & 4  & \texttt{test\_method} (ensemble branch) & Distance-gated blend \\
  & -- & \texttt{\_distance\_llm\_weight} & Blend weight helper added \\
\bottomrule
\end{tabular}
\end{table}

\bigskip
\noindent The net projection is that the Hybrid closes or exceeds the Pure
LLM's $R^2 > 0.99$ rate of $83.9\%$.  The v3.0 benchmark confirms this
projection: HypatiaX achieves \textbf{89.2\%} near-perfect success rate
($R^2>0.99$, $n=74$, zero catastrophic failures), combining the LLM's
structural extrapolation accuracy with the NN's ability to correct
in-distribution calibration errors.

% ====================================================================
% NEW APPENDICES (extracted from jmlr_paper_final_fixed.tex)
% ====================================================================
\appendix

\section{Complete Proof of Finding~\ref{thm:llm_reproductive}}
\label{app:proof}

\begin{proof}[Detailed Proof]
Let $M_\theta$ be an autoregressive language model with parameters $\theta$ trained on corpus $\mathcal{C}$. Expression $E = t_1 \cdots t_n$ is generated via:
\[
P_\theta(E \mid D) = \prod_{i=1}^n P_\theta(t_i \mid t_{<i}, D)
\]

where $t_{<i} = t_1, \ldots, t_{i-1}$ is the prefix and $D$ represents conditioning data.

\paragraph{Step 1: Training concentrates probability mass.}

Maximum likelihood training optimizes:
\[
\theta^* = \arg\max_\theta \mathbb{E}_{E \sim \mathcal{C}}[\log P_\theta(E)]
\]

This causes $P_{\theta^*}(t_i \mid t_{<i}, D)$ to concentrate on tokens frequently observed in $\mathcal{C}$ given context $t_{<i}$.

\paragraph{Step 2: Novel tokens incur exponential penalty.}

For expression $E \notin \mathcal{C}$, define:
\[
d(E, \mathcal{C}) = \min_{E' \in \mathcal{C}} d_{\text{edit}}(E, E')
\]

Let $\mathcal{N}(E) \subseteq \{1, \ldots, n\}$ denote positions where $t_i$ is novel (low probability given $t_{<i}$). Under independence approximation:
\[
P_\theta(E \mid D) \approx \prod_{i \in \mathcal{N}(E)} p_i \cdot \prod_{i \notin \mathcal{N}(E)} p_i'
\]

where $p_i \ll 1$ for novel tokens and $p_i' \approx 1$ for familiar tokens.

\paragraph{Step 3: Exponential decay.}

Since $|\mathcal{N}(E)| \geq d(E, \mathcal{C})$ and $p_i \leq p_{\max} < 1$:
\[
P_\theta(E \mid D) \leq p_{\max}^{d(E, \mathcal{C})} = \exp(-\beta \cdot d(E, \mathcal{C}))
\]

where $\beta = -\log p_{\max} > 0$. Thus generation probability decays exponentially with edit distance from training data, establishing reproductive behavior. \qed
\end{proof}

\paragraph{Remark on rigour:} This proof assumes approximate local independence of token probabilities. A fully rigorous treatment would analyse the joint distribution over all tokens. The exponential decay result is a consequence of the independence approximation; empirical results are consistent with this theoretical prediction, but the approximation means the proof should be understood as providing intuition rather than a formal guarantee\citep{kambhampati2024llms}.

% ============================================================================
% APPENDIX B: HYPERPARAMETER SETTINGS
% ============================================================================


\section{Reproducibility Statement}
\label{app:reproducibility}

All experimental data, source code, and configuration files are publicly
available to ensure full reproducibility of our results.

\subsection{Code Repository}

\textbf{GitHub:} \url{https://github.com/sednabcn/LLM-HypatiaX-Papers}

\paragraph{Repository Structure:}
\begin{small}
\begin{verbatim}
LLM-HypatiaX-Papers/
+-- papers/
|   +-- 2025-JMLR/
|       +-- hypatiax/
|           +-- core/
|           |   +-- base_pure_llm/     # Pure LLM baseline
|           |   +-- generation/        # Hybrid systems
|           |   +-- training/          # Neural network baseline
|           +-- data/results/          # Raw experimental data
|           +-- experiments/
|           |   +-- benchmarks/        # Feynman, noise-sweep runners
|           |   +-- comparison/        # Comparative analysis
|           |   +-- tests/             # DeFi extrapolation tests
|           +-- protocols/             # Experimental protocols (v2)
|           +-- tools/
|               +-- symbolic/          # PySR / hybrid_system_v40
|               +-- validation/        # Multi-layer validators
|       +-- figures/                   # 13 generated figures (PDF/PNG)
|       +-- paper/                     # LaTeX manuscript
|       +-- reports/                   # Analysis reports
|       +-- supplementaries/           # Supplementary materials S1-S8
+-- docs/
+-- templates/
+-- requirements.txt
\end{verbatim}
\end{small}

\subsection{Running Experiments}

\paragraph{Quick Start (all v2-corrected commands):}
\begin{small}
\begin{verbatim}
git clone https://github.com/sednabcn/LLM-HypatiaX-Papers
cd LLM-HypatiaX-Papers/papers/2025-JMLR
pip install -r ../../requirements.txt
source activate_hypatiax.sh

# Core 15 benchmark (§6.4)
python hypatiax/experiments/comparison/\
ultimate_comparative_suite_complete_.py \
    --seed 42 --symbolic-timeout 1800 --v2

# DeFi 74-case benchmark v3.0 (§6.5)
python hypatiax/experiments/tests/\
test_enhanced_defi_extrapolation.py \
    --fixed-denominator 66 --nan-penalty --v2

# Feynman SR — Phase 3 (§5.8, ~1h 13min on reference hardware)
python hypatiax/experiments/benchmarks/\
run_comparative_suite_benchmark_v2.py \
    --noiseless --threshold 0.9999 \
    --nn-seeds 3 --samples 200 \
    --method-timeout 900 --pysr-timeout 900

# Statistical analysis
python hypatiax/analysis/statistical_analysis_full.py --v2

# All 13 figures
python supplementaries/generate_figures/generate_figures.py --v2
\end{verbatim}
\end{small}

\subsection{Hardware and Software}

All experiments were run on a single machine: Intel Celeron T3100 @
1.90\,GHz (2~cores), 3\,GB DDR2 RAM, no GPU, Kali GNU/Linux Rolling.
Runtimes on modern multi-core hardware will be substantially lower; a GPU
is \emph{not} required to reproduce the core results.

\begin{table}[H]
\centering
\caption{Software environment (verified 18 March 2026).}
\label{tab:runtime_reproducibility}
\small
\begin{tabular}{@{}lll@{}}
\toprule
\textbf{Component} & \textbf{Version} & \textbf{Note} \\
\midrule
Python       & 3.12.2  & \\
PySR         & 1.5.9   & \\
Julia        & 1.12.2  & required by PySR \\
juliacall    & 0.9.26  & Python$\leftrightarrow$Julia bridge \\
NumPy        & 2.3.5   & \\
SciPy        & 1.16.3  & \\
SymPy        & 1.14.0  & \\
Pandas       & 2.3.3   & \\
Matplotlib   & 3.10.7  & \\
Seaborn      & 0.13.2  & \\
scikit-learn & 1.7.2   & \\
Anthropic SDK & 0.73.0 & \texttt{claude-sonnet-4-20250514}, temp~0.0 \\
\midrule
PySR populations & 2 & hardware limit (2-core CPU) \\
\bottomrule
\end{tabular}
\end{table}

\paragraph{Expected Wall-Clock Times (reference hardware):}
\begin{small}
\begin{tabular}{@{}lcc@{}}
\toprule
\textbf{Experiment} & \textbf{Tests} & \textbf{Wall Time} \\
\midrule
Pure LLM Baseline          & 40  & 5--10 min \\
Pure Symbolic (Core 15)    & 18  & 117 min \\
LLM-Guided Hybrid          & 30  & 35 min \\
DeFi Suite (73 cases)      & 23  & 44 min \\
Hybrid LLM+NN (v40)        & 30  & 45 min \\
Feynman SR Phase~3         & 30  & $\sim$73 min \\
Statistical Analysis       & --- & 15 min \\
Figure Generation (13)     & --- & 10 min \\
\midrule
\textbf{All Experiments}   & \textbf{131} & \textbf{$\sim$6 hours} \\
\bottomrule
\end{tabular}
\end{small}

\subsection{Random Seeds}

\begin{itemize}
\item All campaigns: base seed \texttt{42} (\texttt{BenchmarkProtocol.seed = 42})
\item Per-equation offset: \texttt{seed + index $\times$ 7}
\item LLM calls: temperature \texttt{0.0} (deterministic) for all methods
\item Neural network (Phase~3): 3 independent seeds per equation; median $R^2$ reported
\end{itemize}

\subsection{Measurement Correction (v1\,$\to$\,v2)}

A bug in \texttt{evaluate\_llm\_formula} was corrected on 2026-03-16.
The function used a hardcoded absolute threshold (\texttt{ss\_tot > 1e-10})
that silently returned $R^2=0$ for equations with outputs far below unity.
The fix scales the threshold relative to the data magnitude.
All results in this paper use the v2 corrected harness.
See Appendix~\ref{app:repro_statement} and Section~\ref{sec:r2_bugfix}
for the exact code change and affected equations.

\subsection{Troubleshooting}

\paragraph{PySR / Julia installation:}
\begin{small}
\begin{verbatim}
# Install Julia 1.12.2 (required)
wget https://julialang.org/downloads/  # follow official installer
export PATH="$HOME/.juliaup/bin:$PATH"

# Install PySR
pip install pysr==1.5.9
python -c "import pysr; pysr.install()"
\end{verbatim}
\end{small}

\paragraph{Arrhenius hang (Feynman test~4):}
Always use \texttt{--method-timeout 900}. Without it, PySR enters an
infinite search on this single-variable flat-landscape equation.
This is a known Julia JIT behaviour on 2-core hardware; counted as a
genuine failure in the paper results.

\paragraph{LLM API:}
\begin{small}
\begin{verbatim}
export ANTHROPIC_API_KEY="your-api-key-here"
\end{verbatim}
\end{small}

\paragraph{Memory on constrained hardware:}
\begin{small}
\begin{verbatim}
# In hypatiax/tools/symbolic/hybrid_system_v40.py
# populations=2  (hardware-limited default used in this paper)
\end{verbatim}
\end{small}

\subsection{Reproducibility Checklist}

\begin{enumerate}
\item Clone repository and install dependencies (including Julia 1.12.2)
\item Set \texttt{ANTHROPIC\_API\_KEY} for LLM-dependent runners
\item Run Core~15 benchmark — confirm Mann-Whitney U=0, p$<10^{-6}$
\item Run DeFi benchmark — confirm HypatiaX 89.2\% vs LLM 62.2\% (n=74, v3.0)
\item Run Feynman Phase~3 — confirm Hybrid DeFi 96.7\% ($R^2>0.9999$)
\item Run statistical analysis — confirm Cohen's d = 0.95 (pooled)
\item Regenerate 13 figures — confirm visual match with paper
\item Verify v2 result files match checksums in repository README
\end{enumerate}

\subsection{Contact}

\begin{itemize}
\item \textbf{GitHub Issues:} \url{https://github.com/sednabcn/LLM-HypatiaX-Papers/issues}
\item \textbf{Email:} \texttt{ruperto.bonet@modelphysmat.com}
\item \textbf{Response time:} Within 2 weeks
\item \textbf{Documentation:} \path{docs/QUICK_START_GUIDE.md}
\end{itemize}

We maintain the repository and welcome contributions, bug reports, and extension requests.

% ============================================================================
% APPENDIX G: COMPUTATIONAL COMPLEXITY ANALYSIS
% ============================================================================


\section{Computational Complexity Analysis}
\label{app:complexity}

\subsection{Theoretical Complexity}

\begin{complexity}[HypatiaX Complexity]
Let $n$ be candidate expressions, $m$ evaluation points, $k$ symbolic constraints. Total complexity is $\mathcal{O}(n \cdot (k + m))$.
\end{complexity}

\begin{proof}
For each candidate expression $E$:
\begin{enumerate}
\item \textbf{Symbolic constraint checking}: $\mathcal{O}(k)$ operations (dimensional analysis, domain rules)
\item \textbf{Numerical evaluation}: $\mathcal{O}(m)$ operations (fitness computation over dataset)
\end{enumerate}

Total per candidate: $\mathcal{O}(k + m)$. With $n$ candidates: $\mathcal{O}(n(k + m))$. \qed
\end{proof}

\paragraph{Comparison with Neural Networks:}
\begin{itemize}
\item \textbf{Neural Network Training}: $\mathcal{O}(e \cdot b \cdot (d \cdot h + h^2))$ where $e$ = epochs, $b$ = batch size, $d$ = input dimension, $h$ = hidden units
\item \textbf{Typical values}: $e=200$, $b=160$, $d=3$, $h=64$ $\Rightarrow$ $\mathcal{O}(2 \times 10^6)$ operations
\item \textbf{HypatiaX}: $n=1000$, $m=160$, $k=50$ $\Rightarrow$ $\mathcal{O}(2 \times 10^5)$ operations per iteration
\end{itemize}

Despite comparable per-iteration complexity, symbolic methods require more iterations (500 vs 200) for convergence, explaining the 27$\times$ runtime difference.

\subsection{Empirical Timing Breakdown}

Profiling HypatiaX on representative test cases:

\begin{table}[H]
\centering
\caption{Computational Cost Breakdown (average per test)}
\label{tab:timing_breakdown}
\begin{tabular}{@{}lcc@{}}
\toprule
\textbf{Component} & \textbf{Time (s)} & \textbf{Percentage} \\
\midrule
\multicolumn{3}{c}{\textit{Pure Symbolic (PySR)}} \\
\midrule
Population initialization & 5.2 & 1.3\% \\
Evolutionary search & 365.8 & 93.9\% \\
Symbolic simplification & 12.3 & 3.2\% \\
Validation framework & 6.7 & 1.6\% \\
\midrule
\textbf{Total} & \textbf{390.0} & \textbf{100\%} \\
\midrule
\multicolumn{3}{c}{\textit{LLM-Guided Hybrid}} \\
\midrule
LLM API call & 8.5 & 12.1\% \\
Expression parsing & 2.0 & 2.9\% \\
Initial validation & 3.0 & 4.3\% \\
PySR refinement (when needed) & 52.5 & 75.0\% \\
Final validation & 4.0 & 5.7\% \\
\midrule
\textbf{Total} & \textbf{70.0} & \textbf{100\%} \\
\midrule
\multicolumn{3}{c}{\textit{Neural Network}} \\
\midrule
Data preparation & 0.3 & 17.6\% \\
Training (200 epochs) & 1.2 & 70.6\% \\
Validation & 0.2 & 11.8\% \\
\midrule
\textbf{Total} & \textbf{1.7} & \textbf{100\%} \\
\bottomrule
\end{tabular}
\end{table}

\paragraph{Key Findings:}
\begin{itemize}
\item Evolutionary search dominates (93.9\%) in pure symbolic mode
\item LLM-guided mode routes 68 of 74 cases (92\,\%) directly to the LLM formula, bypassing full PySR search; median speedup on those cases is $1.73\times$ over the neural baseline \citep{kambhampati2024llms}
\item Validation overhead is negligible (1.6-5.7\%)
\item Neural networks are 27$\times$ faster but achieve 1231\% extrapolation error
\end{itemize}

\subsection{Scalability Analysis}

Testing scalability with varying dataset sizes:

\begin{table}[H]
\centering
\caption{Scalability with Dataset Size}
\label{tab:scalability}
\begin{tabular}{@{}cccc@{}}
\toprule
\textbf{Data Points} & \textbf{Pure PySR (s)} & \textbf{LLM-Hybrid (s)} & \textbf{Speedup} \\
\midrule
100 & 245 & 48 & 5.1$\times$ \\
500 & 390 & 70 & 5.6$\times$ \\
1000 & 521 & 95 & 5.5$\times$ \\
5000 & 1245 & 215 & 5.8$\times$ \\
10000 & 2340 & 410 & 5.7$\times$ \\
\bottomrule
\end{tabular}
\end{table}

\paragraph{Observation:} Speedup remains relatively constant (5.1--5.8$\times$) across dataset sizes, consistent with LLM warm-starting providing efficiency gains independent of data volume\citep{kambhampati2024llms}.

\subsection{Computational Cost vs Accuracy Tradeoff}

\begin{table}[H]
\centering
\caption{Runtime vs Extrapolation Error Tradeoff}
\label{tab:cost_accuracy_tradeoff}
\begin{tabular}{lccc}
\toprule
\textbf{Method} & \textbf{Runtime (s)} & \textbf{Interpolation $R^2$} & \textbf{Extrapolation Error} \\
\midrule
Neural Network & 1.7 & 0.940 & 1231\% \\
LLM-Guided Hybrid & 70.0 & 0.931 & 0.0\% \\
Pure PySR & 390.0 & 0.982 & 0.0\% \\
\midrule
\multicolumn{4}{l}{\textit{Cost per percentage point of extrapolation error reduction:}} \\
\multicolumn{4}{l}{Neural Network $\to$ Hybrid: \textbf{0.056 s per \% error reduced}} \\
\multicolumn{4}{l}{Hybrid $\to$ Pure PySR: \textbf{$\infty$} (both achieve near-zero extrap.\ error on this benchmark)} \\
\bottomrule
\end{tabular}
\end{table}

\paragraph{Interpretation:} The 41$\times$ runtime increase (1.7s $\to$ 70s) reduced mean extrapolation error from 1,231\% (neural network) to near-zero (median $< 10^{-12}$ relative). For scientific applications where extrapolation is critical, this is an acceptable tradeoff. Pure PySR provides no additional benefit over Hybrid despite 5.6$\times$ longer runtime.

\subsection{Parallelisation and Future Speedup}

PySR was configured with 2 parallel populations on the experimental hardware (Intel Celeron T3100, 2-core CPU). On modern multi-core hardware (e.g.\ 16+ cores), PySR can exploit 12 or more parallel populations, yielding an estimated ${\sim}8\times$ wall-clock speedup relative to our reported runtimes. GPU parallelisation of fitness evaluation (75\% of runtime) could further reduce total runtime to ${\sim}1$--$2$ seconds per test, making the hybrid approach competitive with neural networks on speed while retaining near-zero extrapolation error.

% ============================================================================
% APPENDICES H & I: SUPPLEMENTARY MATERIALS AND EXTENDED DISCUSSION
% Corrected values integrated
% Updated: 2026-01-27
% ============================================================================

% ============================================================================
% APPENDIX H: SUPPLEMENTARY MATERIALS
% ============================================================================


\section{Extended Discussion}
\label{app:extended_discussion}

\subsection{Why LLMs Fail on Novel Domains}

Our empirical results (95\% classical vs 40\% DeFi, Table~\ref{tab:llm_detailed}) raise a question: \textbf{Why does LLM performance decline on specialised domains?}\citep{kambhampati2024llms}

\paragraph{Training Data Coverage Hypothesis:}

Classical physics formulas appear extensively in:
\begin{itemize}
\item \textbf{Textbooks}: Digitized and crawled by Common Crawl, Books3 corpus
\item \textbf{Wikipedia}: High-quality articles with LaTeX equation markup
\item \textbf{Academic papers}: Often reproduce fundamental equations in introduction sections
\item \textbf{Code repositories}: Physics simulations, homework solutions (GitHub, GitLab)
\item \textbf{Q\&A forums}: StackExchange Physics, Mathematics, Computational Science
\item \textbf{Educational websites}: Khan Academy, HyperPhysics, Physics Classroom
\end{itemize}

\textbf{Estimated training exposure}: $10^6$-$10^7$ occurrences for common equations like $F = ma$, $E = mc^2$, $PV = nRT$.

DeFi formulas, by contrast:
\begin{itemize}
\item \textbf{Emerged post-2018}: After many LLM training data cutoffs (GPT-3: Sept 2021)\citep{kambhampati2024llms}
\item \textbf{Limited sources}: Primarily white papers, technical docs, GitHub repos
\item \textbf{Non-standard notation}: $x \cdot y = k$ vs traditional mathematical conventions
\item \textbf{Niche audience}: Cryptocurrency/blockchain community, not mainstream education
\item \textbf{Rapid evolution}: Formulas change as protocols update (Uniswap v2 $\to$ v3 $\to$ v4)
\end{itemize}

\textbf{Estimated training exposure}: $10^2$-$10^3$ occurrences for common DeFi equations.

\paragraph{Empirical Validation of Coverage Hypothesis:}

We tested this by searching for equation occurrences in Common Crawl snapshots:

\begin{table}[H]
\centering
\caption{Equation Prevalence in Training Data (estimated)}
\label{tab:equation_prevalence}
\begin{tabular}{@{}lcc@{}}
\toprule
\textbf{Equation} & \textbf{Domain} & \textbf{Approx. Occurrences} \\
\midrule
$F = ma$ & Classical Physics & $\sim 5 \times 10^6$ \\
$E = \frac{1}{2}mv^2$ & Classical Physics & $\sim 2 \times 10^6$ \\
$PV = nRT$ & Classical Physics & $\sim 1 \times 10^6$ \\
$k = A\exp(-E_a/RT)$ & Chemistry & $\sim 5 \times 10^4$ \\
pH = pKa + $\log([A^{-}]/[HA])$ & Chemistry & $\sim 2 \times 10^4$ \\
$x \cdot y = k$ & DeFi & $\sim 3 \times 10^2$ \\
IL = 2$\sqrt{r}/(1+r)$ - 1 & DeFi & $\sim 50$ \\
VaR = P$\sigma$z & Finance (general) & $\sim 1 \times 10^4$ \\
\bottomrule
\end{tabular}
\end{table}

\textbf{Correlation}: LLM success rate vs log(occurrences): Spearman $\rho = 0.89$, $p < 0.001$.

\paragraph{Conceptual Complexity Comparison:}

Beyond data availability, specialized domains require distinct reasoning patterns:

\begin{table}[H]
\centering
\caption{Conceptual Requirements Comparison}
\label{tab:conceptual_complexity}
\begin{tabular}{@{}lcc@{}}
\toprule
\textbf{Requirement} & \textbf{Classical Physics} & \textbf{DeFi/Risk} \\
\midrule
Distributional reasoning & Rare & Essential \\
Multi-step derivations & Common but formulaic & Domain-specific \\
Statistical API knowledge & Basic (mean, std) & Advanced (scipy.stats) \\
Convention awareness & Well-established & Evolving \\
Asymptotic analysis & Standard & Context-dependent \\
Numerical edge cases & Predictable & Subtle (e.g., VaR tails) \\
Unit consistency & Strict & Often dimensionless \\
\bottomrule
\end{tabular}
\end{table}

\paragraph{Example: VaR Calculation Failure}

Pure LLM generated for 95\% parametric VaR:
\begin{verbatim}
VaR_95 = portfolio_value * sigma * (-1.645)  # WRONG SIGN
\end{verbatim}

\textbf{Failure mode}: LLM confused sign convention. In finance, VaR is reported as positive number representing potential loss, but calculation uses negative z-score. This subtle convention is not well-represented in training data\citep{lewkowycz2022minerva}.

\textbf{Result}: $R^2 = -10.05$ (predictions anticorrelated with ground truth)

\textbf{Symbolic rescue}: PySR discovered correct relationship including sign, achieving $R^2 = 0.9988$.

\subsection{Future Research Directions}
\label{sec:future_work}

\paragraph{Near-Term (1-2 years):}

\begin{enumerate}
\item \textbf{GPU-Accelerated Symbolic Regression}:
   \begin{itemize}
   \item Parallelize fitness evaluation across thousands of GPU cores
   \item GPU acceleration may offer substantial speedup over current CPU-based search
   \item Would potentially make hybrid approach more competitive with neural networks on speed
   \end{itemize}

\item \textbf{Active Learning Integration}:
   \begin{itemize}
   \item Use LLM feedback to guide symbolic search
   \item Potentially reduce required data through uncertainty-guided sampling
   \item Particularly relevant for expensive experimental data
   \end{itemize}

\item \textbf{Multi-Scale Constant Handling}:
   \begin{itemize}
   \item Logarithmic preprocessing for constants spanning >6 orders of magnitude
   \item Adaptive search ranges based on data scale
   \item May address the gravitational force failure observed in current experiments (1/15 failure rate)
   \end{itemize}
\end{enumerate}

\paragraph{Medium-Term (3-5 years):}

\begin{enumerate}
\item \textbf{PDE Discovery}:
   \begin{itemize}
   \item Combine Fourier Neural Operators with symbolic regression
   \item Discover spatial/temporal derivative structures
   \item Applications: fluid dynamics, heat transfer, wave propagation
   \end{itemize}

\item \textbf{Hierarchical Discovery}:
   \begin{itemize}
   \item Discover sub-expressions separately, then combine
   \item Example: $f(x) = g(h(x))$ $\to$ first find $h$, then $g$
   \item Reduces search space exponentially
   \end{itemize}

\item \textbf{Transfer Learning}:
   \begin{itemize}
   \item Leverage formulas from related domains
   \item Chemical kinetics $\to$ biochemical pathways
   \item Classical mechanics $\to$ robotics control
   \end{itemize}
\end{enumerate}

\paragraph{Long-Term (5-10 years):}

\begin{enumerate}
\item \textbf{Automated Theory Formation}:
   \begin{itemize}
   \item Long-term aspiration: not just discovering equations, but underlying principles
   \item Example: Discovering conservation laws from data
   \item Would require integration with automated theorem proving
   \end{itemize}

\item \textbf{Multi-Modal Scientific Discovery}:
   \begin{itemize}
   \item Incorporate experimental images, diagrams, protocols
   \item Vision-language models for lab automation
   \item End-to-end hypothesis $\to$ experiment $\to$ analysis loops
   \end{itemize}

\item \textbf{Uncertainty-Aware Discovery}:
   \begin{itemize}
   \item Discover not just $f(x)$, but $f(x) \pm \sigma(x)$
   \item Bayesian symbolic regression
   \item Applications: risk-sensitive decision making
   \end{itemize}
\end{enumerate}

\subsection{Final Thoughts}

The central lesson of this work is that \textbf{retrieval of familiar patterns may differ from structural discovery}\citep{kambhampati2024llms}. On the evaluated benchmark, LLMs achieved high success rates on widely-documented equations (approximately 95\% on classical physics) while showing substantially lower reliability in specialized domains requiring genuine structural inference (approximately 40\% on DeFi; extrapolation error statistics in Table~\ref{tab:five_systems_full}).

This distinction matters for science. Discovery requires:
\begin{enumerate}
\item Extrapolation to novel regimes (symbolic near-zero vs substantially larger neural error; Table~\ref{tab:five_systems_full})
\item Functional form recovery (interpretability, insight)
\item Reliable predictions for decision-making (drug dosing, climate policy)
\end{enumerate}

LLM-only approaches showed inconsistency on all three for specialized domains on the evaluated benchmark. Hybrid architectures combining LLM interfaces with symbolic engines addressed all three, at the cost of approximately 27$\times$ longer runtime---an acceptable tradeoff when correctness matters.

As we integrate AI more deeply into scientific workflows, maintaining this distinction between approximation and discovery will be increasingly important. A productive direction is the orchestration of complementary capabilities: LLMs for communication and initialisation, symbolic methods for structured mathematical search, and human experts for interpretation and validation.



% =============================================================================
% INSERTED FROM jmlr-hypatiax-paper-final.tex — Appendix B & C
% =============================================================================



\bibliographystyle{plainnat}
\bibliography{references}
\end{document}
